\section{Preliminaries}\label{sec:prelim}

\subsection{Machines and the statement}\label{sec:prelim-machines}

We use the model and the formal statement of the first paper~\cite[Section~2]{FirstPaper2026}. A machine has a
read-only input tape, $k$ work tapes and a write-only output tape. A $d\times d$ matrix $A$ with entries of
$b=b(d)=\max(1,\lceil\log_2d\rceil)$ bits is encoded row by row as a word of $L=d^2b$ bits, and the machine must append
the encoding of the transpose. Statement~(T) is the Lean proposition \lean{Transpose.LinearTimeTransposable}: some
machine and some constant $c$ do this within $c(L+1)$ steps for all $d\ge1$ and all~$A$. A machine is
\emph{oblivious} at size~$d$ if the positions of all heads and the times at which output bits are appended are the
same for all inputs of that size.

\subsection{Whole-row codes}\label{sec:prelim-codes}

Rows and columns of a $d\times d$ array are indexed by $[d]=\{0,\dots,d-1\}$, and an array over an alphabet~$\Sigma$
is an element of $\Sigma^{[d]\times[d]}$. A \emph{helper map} is a map $\rho$ from $[d]$ to the subsets of $[d]$; its
\emph{window} is $\max_j\card{\rho(j)}$. We read $\rho(j)$ as the set of rows that help column~$j$.

\begin{definition}[Whole-row codes]\label{def:wholerow}
\formal{StatementsRowSpread.lean}{RowSpread.wholeRowHelpers, RowSpread.IsWholeRowCodeExcept, RowSpread.WholeRowCardinalityBound, RowSpread.UniformWholeRowRateCanonicalBits}
A set $C\subseteq\Sigma^{[d]\times[d]}$ is a \emph{whole-row code for~$\rho$} if for every $j\in[d]$ any two arrays of
$C$ that agree in all entries of the rows $\rho(j)$ agree in all entries of column~$j$. It is \emph{two-sided} if in
addition any two arrays of $C$ that agree in all entries of the columns $\rho(j)$ agree in all entries of row~$j$.
\end{definition}

The function that computes a column from its helper rows is arbitrary, and differs from column to column. An entry
$(i,j)$ with $i\in\rho(j)$ lies in one of its own helper rows; we call it \emph{incident}. There are at most $wd$
incident entries, so they do not affect rates. These are the $\rho$-codes of the
synthesis~\cite[Section~6.7]{SeventeenthPaper2026}.

We also read a helper map as the relation $\{(i,j):i\in\rho(j)\}$ on $[d]$, and as the bipartite graph that joins
row~$i$ to column~$j$ when $i\in\rho(j)$. The relation is \emph{symmetric} if $i\in\rho(j)$ exactly when
$j\in\rho(i)$, and \emph{reflexive} if $j\in\rho(j)$ for all~$j$. Its \emph{incidence girth} is the girth of this
bipartite graph. An \emph{empty $r\times r$ incidence rectangle} is a set $I$ of $r$ rows together with a set $J$ of
$r$ columns such that $i\notin\rho(j)$ for all $i\in I$ and $j\in J$.

\begin{definition}[Rates]\label{def:rate}
The \emph{rate} of a nonempty code $C$ over $\Sigma$ is $\log_{\card\Sigma}\card C/d^2$. For $w\ge1$ let $R(w)$ be the
upper limit, as $d\to\infty$, of the largest rate of a whole-row code of window at most $w$ on $d\times d$ arrays with
entries in $\{0,1\}^{b(d)}$, the \emph{canonical alphabet}.
\end{definition}

\begin{example}[Pairing]\label{ex:pairing}
For $\rho(j)=\{j\}$ the symmetric arrays form a two-sided whole-row code of rate $\tfrac12+\tfrac1{2d}$. So
$R(w)\ge\tfrac12$ for every~$w$.
\end{example}

In this notation \cref{conj:MCC} says $\sup_wR(w)<1$; its Lean form is
\lean{Wall.CommonDeficit}. The Lean form bounds the size by $2^{(1-\delta+\eta)L}$ for every $\eta>0$ and all large~$d$;
this is equivalent to the form of \cref{conj:MCC}, with $\delta$ halved in one direction.

\begin{conjecture}[Half-rate conjecture]\label{conj:half}
\formal{StatementsRowSpread.lean}{RowSpread.WholeRowHalfRateCanonicalBits}
$R(w)=\tfrac12$ for every $w\ge1$.
\end{conjecture}

\begin{proposition}[Known bounds]\label{prop:known}
\partformal{KnownBounds.lean}{KnownBounds.cardinality_pow, KnownBounds.cardinality_bound, KnownBounds.window_one, KnownBounds.cardinality_pow_threshold, KnownBounds.uniform_in_alphabet}{The sharper window-two bound remains the written proof of \cite[Proposition~6.22(b)]{SeventeenthPaper2026}.}
For every alphabet, a whole-row code of window $w$ has rate at most $\tfrac{w}{w+1}(1+\tfrac1d)$, and for $w\le2$ it
has at most $\card\Sigma^{d(d+1)/2+d-1}$ elements. In particular $R(1)=R(2)=\tfrac12$ and
$\tfrac12\le R(w)\le\tfrac{w}{w+1}$.
\end{proposition}

\begin{outline}
The averaging argument is that of \cite[Proposition~6.24(d)]{SeventeenthPaper2026}. Fix a row order and retain, in
each column, the entries up to its last helper row. These entries determine the array by induction in that order.
For a set of $k$ helper rows and one other row, each of the $k+1$ rows is last in equally many orders: exchanging
two row labels is a bijection between the two classes of orders. Summing these counts gives an order of cost at most
$wd(d+1)/(w+1)$. The Lean proof retains the sharper numerator $wd^2+\sum_j\card{\rho(j)}$, and applies to rectangular
arrays as well. In particular
$\card C^{w+1}\le\card\Sigma^{wd(d+1)}$, with no linearity or restriction on the recovery functions.
At window one this gives $\card C\le\card\Sigma^{d(d+1)/2}$; pairing attains it. The threshold for the limiting
upper rate is independent of the alphabet. Window two uses the separate multigraph-ranking argument cited above.
\end{outline}

\begin{remark}[Windows that grow with $d$]\label{rem:growing}
\partformal{GraphMatchingCode.lean}{GraphMatchingCode.exists_code, GraphMatchingCode.mem_code_iff, GraphMatchingCode.cardinality, GraphMatchingCode.isWholeRowCode, GraphMatchingCode.window_iff, GraphMatchingCode.update_nonedge, GraphMatchingCode.mem_linearCode_iff_code, GraphMatchingCode.finrank_linearCode_add}{The graph-existence results and their asymptotic parameters, cited below. The finite graph-to-code conversion is formal.}
In \cref{def:rate,conj:MCC} the window is fixed before $d$ grows. This cannot be relaxed much. Mittal and Raz
conjectured that for some constant $c>0$ no set of $2^{(1-c)d^2}$ matrices of bits has every column a function of
$cd$ rows; in their terms, that the transpose is a block-rigid function~\cite[Conjecture~36]{MittalRaz2021}. In the
terms of this paper that is a deficit $c$ at window $cd$, for entries of one bit. It is false. By a construction of
Braverman and Forbes, given in Mittal's dissertation, for every $\eps>0$ there is an $\alpha>0$ such that for
infinitely many~$d$ some helper map of window at most $d^{1-\alpha}$ has an $\F_2$-linear whole-row code of
codimension at most $d^{1+\eps}$~\cite[Question~10.8.2, Proposition~10.8.3 and Lemma~10.8.5]{Mittal2025}. Used for
each bit of the entries, it gives whole-row codes over the canonical alphabet of rate at least $1-d^{\eps-1}$. The
codes are defined by parities over the induced matchings of a dense Ruzsa--Szemer\'edi graph, and the same
constructions are delivery schemes in coded caching~\cite{ShanmugamEtAl2017,YanEtAl2017}; see the second paper of the
series~\cite{SecondPaper2026}. So the best rate at window $d^{1-\alpha}$ tends to one, and \cref{conj:MCC} would be
false for windows of that size. A second construction of Alon, Moitra and Sudakov lowers the window. For
infinitely many~$d$ it is a graph on $d$ vertices, the words of a fixed length with an edge between two words that
agree in few positions, that misses at most $d^{3/2+\eps}$ edges and whose edges are partitioned into
$d^{2-c\eps^3}$ induced matchings~\cite[Claim~2.11, Lemma~2.12, Theorem~2.20 and the paragraph after it]{AMS2013}. All its
vertices miss equally many others, hence at most $2d^{1/2+\eps}$. Take $\rho(j)$ to be $j$ together with the vertices
missed by~$j$; listing each edge of a matching in both orders gives a partition of the non-incident entries as in
\cref{prop:matchings}, and with it whole-row codes of rate at least $1-2d^{-c\eps^3}$ at window at most
$2d^{1/2+\eps}+1$. At a fixed window, codes of this form have rate at most one half in the limit
(\cref{prop:matchings}).
\end{remark}

\begin{outline}[The finite conversion in \cref{rem:growing}]
Let $G$ be a simple graph on $d$ vertices whose edges are partitioned into $m$ nonempty induced matchings.
Set $\rho(j)=\{i:\{i,j\}\notin E(G)\}$, including $j$ itself. Each matching supplies one block of ordered
entries, with both orders of each edge present. Distinct entries in a block are mutual helpers: if their
cross endpoints were adjacent, the induced-subgraph property and the unique neighbour in a matching would
force the two ordered entries to coincide. The blocks therefore form exactly the partition required in
\cref{prop:matchings}.

For any finite abelian group $A$, impose one prescribed sum per block. The resulting whole-row code has
exactly $\card A^{d^2-m}$ words; diagonal and non-edge entries are free. Its helper width at column $j$ is
$d-\deg_G(j)$, so minimum degree at least $d-w$ gives window at most $w$. Over any field the zero-sum code
is a linear subspace of codimension exactly $m$, since the disjoint nonempty checks give a surjective linear
map to the $m$ check values. These are the finite statements used with the cited graph families; no
existence theorem for those families is part of this formalization.
\end{outline}

\subsection{The Tape Coding Conjecture}\label{sec:prelim-tcc}

We recall the objects of the first paper in the form in which \cref{sec:tape} uses
them~\cite[Sections~8 and~9]{FirstPaper2026}. A \emph{time--space diagram} with $k$ work tapes and length~$T$ records,
for every step, the positions of the input head and of the work heads and the length of the output so far; it is
\emph{valid} if every head moves by at most one cell per step. Fix a granularity $G\ge1$ and let $\ell=\lceil db/G\rceil$, the length of a block. The \emph{block network} of
the diagram has as nodes the \emph{slots}, the consecutive stretches of $\ell$ steps, in their order in time; a
\emph{memory edge} from every slot to the next; and a \emph{chord edge} from a slot to a later one whenever some work
head, during the later slot, is in blocks in which it last resided during the earlier one. The rows of the input that
the input head touches during a slot are \emph{attached} to it, and an output bit is produced at the slot in which it
is appended. A memory edge can carry the blocks in which the heads reside and a chord edge one block for every tape
that has the chord; these are the \emph{profile capacities}~\cite[Definitions~8.2 and~8.3]{FirstPaper2026}.
A \emph{causal code} assigns to every
edge a function of the values on the edges into its tail and of the bits of the rows attached there, and to every
output bit a function of the same data at its slot. It \emph{solves complete exchange} on a set of inputs if on those
inputs every edge value fits its capacity and every output bit is the transposed input bit.

\begin{conjecture}[Tape Coding Conjecture, robust form {\cite[Conjecture~9.1]{FirstPaper2026}}]\label{conj:TCC}
For all $k$, $G\ge1$, $c$ and $C_0$ and all sufficiently large~$d$, no causal code solves complete exchange at the
profile capacities on the block network of a valid diagram with $k$ work tapes and length at most $cL$, for a set of
inputs of co-entropy at most $C_0\,d\,b(d)$.
\end{conjecture}

The Lean proposition is \lean{Transpose.Tape.RobustNoCode_TM}. The \emph{plain form}, \lean{NoCode_TM}, asks the same
for codes that are correct on all inputs. The robust form implies that (T)
fails~\cite[Theorem~9.2]{FirstPaper2026}, and the plain form excludes oblivious linear-time
transposers~\cite[Theorem~9.9]{FirstPaper2026}; both implications are Lean-verified there.

\subsection{Notation}\label{sec:prelim-notation}

The letters of the earlier papers collide, so we fix them here.
\begin{center}
\begin{tabular}{@{}ll@{}}
\toprule
$d$, $b=b(d)$, $L=d^2b$ & side length, bits of an entry, bits of an array\\
$\rho$, $w$ & helper map and its window (in the fourth paper $\rho$ is a rate)\\
$U$, $R$ & normalized window; reverse-degree cutoff\\
$\delta$ & deficit: a bound $2^{(1-\delta)L}$ on the size of a code\\
$k$, $c$, $G$ & work tapes, time constant, blocks per row\\
$\ell=\lceil db/G\rceil$ & cells of a block (the $r$ of the first paper)\\
$\theta$, $D$ & density of a cut among the nodes; depth left by the cut\\
$W$ & number of rows that a cone can meet after the cut\\
\bottomrule
\end{tabular}
\end{center}
Logarithms are to base two unless a base is written. A statement holds \emph{for large $d$} if it holds for all $d$
above a threshold that may depend on every other parameter in scope.
